\section{Simultaneous finite models}\label{sec:finite}
\begin{theorem}\label{thm:finite}
Let $h(e)$ be the recognizer's halting time, or $\infty$ if it does not halt. For every even integer $N\ge2$, define
\[
 K_N=\{e<N/2:h(e)<2^N\},
 \qquad\theta_N=\sum_{e\in K_N}4^{-(e+1)}.
\]
There is a flat model of $\mathcal P_K$ on dimension
\begin{equation}\label{eq:finite-dimension}
 d_N=2^{3N+k_N},\qquad k_N=\theta_NN+O(\log(N+2)),
\end{equation}
with every defining defect at most $3\cdot4^{-N}$ in operator norm. The coefficients $\theta_N$ increase to $\thetaK$.
\end{theorem}
\begin{proof}
Take $N$-bit address, fuel and time registers, with the finite coordinate~\eqref{eq:finite-clock}. Their bit zero is least significant. Let the address integer be $x\in\Z/2^N\Z$ and put
\begin{equation}\label{eq:mask}
 f_N(x)=\sum_{e\in K_N}1_{\{x\equiv4^e\pmod{4^{e+1}}\}}.
\end{equation}
The residue classes are disjoint and their periods divide $2^N$, so this is a Boolean function of mean $\theta_N$.

Every interval of $N$ consecutive addresses has
\begin{equation}\label{eq:window}
 \sum_{j=1}^N f_N(x-j)=N\theta_N+O(\log(N+2))
\end{equation}
uniformly in $x$, including cyclic intervals. To prove it, put $m=\lfloor\log_4N\rfloor$. For each $e\le m$, one residue class has a count differing from its mean by at most one. Every flagged address from $e>m$ is a multiple of $4^{m+1}$, so the whole remaining family contributes at most one address to such an interval. Its expected count is at most $N/(3\cdot4^{m+1})$. Summing proves~\eqref{eq:window}.

Set
\begin{equation}\label{eq:slots}
 k_N=1+\max_x\sum_{j=1}^N f_N(x-j)
\end{equation}
and add $k_N$ optional qubits. Let $R_0$ rotate their contents from slot $i$ to slot $i+1$ cyclically. Define
\begin{equation}\label{eq:finite-V}
 V|x,y\rangle=|x+1,R_0^{f_N(x)}y\rangle,
 \qquad F=f_N(x),\quad Q=F P_{\mathrm{slot}\,0},\quad
 Y=(I-F)+FX_{\mathrm{slot}\,0}.
\end{equation}
Define the translates $F_j,Q_j$ by~\eqref{eq:optional-translates} and set
\begin{equation}\label{eq:series}
 B=\sum_{j\ge0}3\cdot4^{-j-1}Q_j,\qquad
 C=\sum_{j\ge0}3\cdot4^{-j-1}F_j.
\end{equation}
These are explicit finite diagonal matrices. Indeed $V$ is a permutation of a finite basis, so the series are periodic geometric sums. They are positive contractions. Their terms commute, so~\eqref{eq:conditional-comm} is exact. Shifting the sums proves~\eqref{eq:conditional-affine} exactly; the controlled flip gives~\eqref{eq:conditional-flip} exactly.

At an active $Q_j$ and a fixed current address $x$, its occupied slot is
\[
 \sum_{r=x-j}^{x-1}f_N(r)\pmod{k_N}.
\]
If $1\le j\le N$, this sum is positive because $f_N(x-j)=1$, and is at most $k_N-1$ by~\eqref{eq:slots}. Thus $Y$ preserves $Q_j$ for those $j$. It always commutes with $C$, and the uncontrolled remainder satisfies
\begin{equation}\label{eq:tail-bound}
 \opn{[Y,VBV^*]}\le2\sum_{j>N}3\cdot4^{-j}=2\cdot4^{-N}.
\end{equation}
This verifies~\eqref{eq:conditional-tail}.

To couple $V$ to the address clock, let $T_0,X_0$ be its ordinary cyclic rotation and first-bit flip and put
\begin{equation}\label{eq:dressing}
 h_N(x)=\sum_{0\le r<x}f_N(r),\quad
 D=\bigoplus_xR_0^{h_N(x)},\quad
 T_a=D(T_0\otimes I)D^*,\quad X_a=D(X_0\otimes I)D^*.
\end{equation}
The coordinate $A_a$ and digit $P_a$ commute with $D$, so all address clock defect bounds from Theorem~\ref{thm:clock-profile} are unchanged. In the $D^*$ gauge, $V$ is ordinary increment except at the all-ones address, where it carries a fixed optional-register permutation. The recursive increment agrees away from that address. At all ones its two output addresses are distinct, with the same optional permutation. The guard is $4^{-N}$ there. Hence its operator defect is $\sqrt2\,4^{-N}$ and its flat defect is $\sqrt2\,4^{-N}2^{-N/2}$.

For $H(x)$ use $4^{-k}$ when the address has $k<N$ initial ones and then a zero, and use zero at all ones. Rotation removes one initial one, so~\eqref{eq:carry-rec} is exact, including its all-ones case. All commutations in~\eqref{eq:carry-step} are exact, and off all ones
\[
 A_a(x+1)-A_a(x)=-1+\tfrac74\,4^{-k}.
\]
At all ones the defect of the last equation is $4^{-N}$. Its flat norm is $4^{-N}2^{-N/2}$. Let $U$ be ordinary cyclic time increment and use ordinary clocks on fuel and time. Their defects have the same bounds as before.

It remains to verify the whole program on these same registers. For a nonzero address let $e$ be the position of its first nonzero base-four digit. Set $R(x)=2^{-e}$ and $L(x)=1$ precisely when that digit is one. At the all-zero address set $R=L=0$. Since $N$ is even, rotation by two bits removes one whole base-four digit. Thus~\eqref{eq:decode-L}--\eqref{eq:decode-R} hold exactly, including at the all-zero address.

For fuel with $k$ initial ones then a zero, put
\[
 Z(q,t)=1_{\{t\equiv0\pmod{2^k}\}}.
\]
On all-ones fuel use $k=N$. When the first fuel bit is zero, $Z=1$. Otherwise the factor $I-P_t$ in~\eqref{eq:reset} requires the first time bit to be zero and the rotated equation tests the remaining $k-1$ bits. On all-ones fuel, the recursive test is true precisely at the all-zero time word. This proves~\eqref{eq:reset} in every case.

At a basis point $(x,q,t)$ assign every workspace counter and control the deterministic configuration
\begin{equation}\label{eq:history}
 \Step^{t\bmod 2^k}\bigl(\Init(x)\bigr).
\end{equation}
These are diagonal entries on existing registers; they add no quantum tensor factor. For ordinary counter entries set $z_i=1,s_i=0$ when $r_i=1$ and $z_i=0,s_i=2r_i$ otherwise; the idle $r_i=0$ case uses $z_i=s_i=0$. These are contractions and give~\eqref{eq:zero-test} exactly. The next time point is initialized if its residue modulo $2^k$ is zero, and is otherwise one step of the old configuration. This is exactly~\eqref{eq:initialize}--\eqref{eq:transition}, also at the cyclic time boundary. A displayed halt occurs in fewer than $2^N$ steps on an input $e<N/2$, so its address is flagged in~\eqref{eq:mask}. Hence~\eqref{eq:halt-flag} is exact.

All workspace and reset variables are diagonal in $(x,q,t)$, and the optional register carries no workspace labels. Every specified remaining commutation therefore holds. The only nonzero defining defects are the clock affine and tail-flip defects, the guarded increments, the carry step and~\eqref{eq:tail-bound}. Their operator norms are bounded by the theorem's constant. The three compulsory registers and the optional register have dimension $2^{3N+k_N}$, and~\eqref{eq:window} gives~\eqref{eq:finite-dimension}. Finally each halting input eventually meets both cutoffs in $K_N$, while the address weights have a summable geometric tail. This proves $\theta_N\uparrow\thetaK$.
\end{proof}
